\chapter{Документация COMSOL-части дипломной работы: численное моделирование термоупругих процессов и экспорт расчётных данных}

\section{Назначение COMSOL-части в общей дипломной работе}

Данная глава описывает COMSOL-часть общей дипломной работы. В рамках проекта COMSOL Multiphysics используется как инженерная среда для построения численной модели, проведения физических симуляций и получения расчётных данных по термоупругому процессу в геологической среде.

COMSOL-часть является базовой инженерной частью проекта, поскольку именно в ней создаётся физическая модель исследуемого процесса. На этом этапе задаются геометрия расчётной области, материалы, физические интерфейсы, начальные и граничные условия, конечно-элементная сетка и параметры временного расчёта. В результате формируется набор численных данных, описывающих изменение температуры, перемещений, скоростей, деформаций, напряжений и свойств материалов во времени и пространстве.

В общей дипломной работе COMSOL-модель выполняет роль источника физических расчётных данных. Эти данные далее могут использоваться в других частях проекта, однако в данной главе основное внимание уделяется именно тому, как была организована COMSOL-симуляция, какие физические процессы в неё заложены и какие данные были получены на выходе.

Основные задачи COMSOL-части:

\begin{enumerate}
    \item построить трёхмерную модель геологической среды и нагретого металлического стержня;
    \item задать свойства материалов для породы и стержня;
    \item настроить теплопередачу в твёрдых телах;
    \item настроить механику твёрдого тела;
    \item связать тепловую и механическую задачи через тепловое расширение;
    \item задать начальные температуры и механические начальные условия;
    \item определить тепловые и механические граничные условия;
    \item построить конечно-элементную сетку;
    \item выполнить нестационарный расчёт во времени;
    \item экспортировать расчётные данные в набор CSV-файлов;
    \item сохранить файл конечно-элементной сетки в формате \texttt{mphtxt};
    \item задокументировать структуру полученного набора данных.
\end{enumerate}

Общая логика COMSOL-части может быть представлена следующим образом:

\begin{equation}
\text{геометрия} \rightarrow \text{материалы} \rightarrow \text{физика} \rightarrow \text{граничные условия} \rightarrow \text{сетка} \rightarrow \text{расчёт} \rightarrow \text{экспорт данных}.
\end{equation}

Таким образом, данная глава является документацией выполненной COMSOL-модели и описывает, какие физические процессы были заложены в симуляцию и какие расчётные данные были получены.

\section{Описание моделируемого физического процесса}

В COMSOL рассматривается задача распространения термоупругого воздействия в геологической среде. Геологическая среда представлена как твёрдое тело, внутри или рядом с которым расположен металлический стержень. Стержень имеет повышенную начальную температуру и выступает источником теплового воздействия.

Из-за разности температур между стержнем и породой возникает тепловой поток. Тепло начинает распространяться от стержня в окружающую геологическую среду. Температурное поле является функцией координат и времени:

\begin{equation}
T = T(x,y,z,t),
\end{equation}

где $x$, $y$, $z$ --- пространственные координаты, а $t$ --- время.

Изменение температуры приводит к тепловому расширению материала. Если разные области среды нагреваются неравномерно, то они расширяются с разной интенсивностью. Это приводит к возникновению перемещений, деформаций и напряжений. Поэтому модель является связанной термоупругой задачей: тепловое поле влияет на механическое состояние среды.

Механический отклик описывается вектором перемещений:

\begin{equation}
\mathbf{u}(x,y,z,t)=
\begin{bmatrix}
u_x(x,y,z,t)\\
u_y(x,y,z,t)\\
u_z(x,y,z,t)
\end{bmatrix}.
\end{equation}

Также в расчёте определяются деформации:

\begin{equation}
\boldsymbol{\varepsilon}=\boldsymbol{\varepsilon}(x,y,z,t),
\end{equation}

и напряжения:

\begin{equation}
\boldsymbol{\sigma}=\boldsymbol{\sigma}(x,y,z,t).
\end{equation}

Итоговая COMSOL-симуляция описывает не только распределение температуры, но и механическую реакцию геологической среды на нагрев.

\section{Математическая основа тепловой части}

Тепловая часть модели описывает распространение температуры в стержне и породе. Основным законом теплопроводности является закон Фурье:

\begin{equation}
\mathbf{q}=-k\nabla T,
\end{equation}

где $\mathbf{q}$ --- вектор плотности теплового потока, $k$ --- коэффициент теплопроводности материала, $\nabla T$ --- градиент температуры.

Знак минус показывает, что тепло распространяется в направлении уменьшения температуры. Если стержень нагрет сильнее, чем порода, тепловой поток направлен от стержня в породу.

Нестационарное уравнение теплопроводности имеет вид:

\begin{equation}
\rho C_p\frac{\partial T}{\partial t}=\nabla\cdot(k\nabla T)+Q,
\end{equation}

где $\rho$ --- плотность материала, $C_p$ --- удельная теплоёмкость, $Q$ --- объёмный источник тепла.

Если теплопроводность материала считается изотропной и постоянной, уравнение можно записать как:

\begin{equation}
\rho C_p\frac{\partial T}{\partial t}=k\nabla^2T+Q.
\end{equation}

Лапласиан температуры в трёхмерной системе координат:

\begin{equation}
\nabla^2T=\frac{\partial^2T}{\partial x^2}+\frac{\partial^2T}{\partial y^2}+\frac{\partial^2T}{\partial z^2}.
\end{equation}

Параметры $\rho$, $C_p$ и $k$ определяют характер распространения тепла. Плотность и теплоёмкость задают тепловую инерционность материала, а теплопроводность определяет скорость распространения тепла.

В данной модели нагретый стержень задаётся как отдельный домен с повышенной начальной температурой. Геологическая среда имеет собственную начальную температуру. После запуска расчёта тепло передаётся от стержня к породе через контактную область.

На границе между стержнем и породой физически должны выполняться условия согласования температуры и теплового потока:

\begin{equation}
T_{rod}=T_{rock},
\end{equation}

\begin{equation}
-k_{rod}\frac{\partial T_{rod}}{\partial n}=-k_{rock}\frac{\partial T_{rock}}{\partial n}.
\end{equation}

Это означает, что температура на общей границе согласована, а тепловой поток, выходящий из стержня, входит в геологическую среду.

\section{Математическая основа механической части}

Механическая часть модели описывает перемещения, деформации и напряжения, возникающие в геологической среде из-за теплового расширения.

Основной неизвестной величиной является вектор перемещений:

\begin{equation}
\mathbf{u}=\begin{bmatrix}u_x\\u_y\\u_z\end{bmatrix}.
\end{equation}

Деформации связаны с перемещениями выражением:

\begin{equation}
\boldsymbol{\varepsilon}=\frac{1}{2}\left(\nabla\mathbf{u}+(\nabla\mathbf{u})^T\right).
\end{equation}

Нормальные компоненты деформации:

\begin{equation}
\varepsilon_{xx}=\frac{\partial u_x}{\partial x},
\end{equation}

\begin{equation}
\varepsilon_{yy}=\frac{\partial u_y}{\partial y},
\end{equation}

\begin{equation}
\varepsilon_{zz}=\frac{\partial u_z}{\partial z}.
\end{equation}

Сдвиговые компоненты деформации:

\begin{equation}
\varepsilon_{xy}=\frac{1}{2}\left(\frac{\partial u_x}{\partial y}+\frac{\partial u_y}{\partial x}\right),
\end{equation}

\begin{equation}
\varepsilon_{yz}=\frac{1}{2}\left(\frac{\partial u_y}{\partial z}+\frac{\partial u_z}{\partial y}\right),
\end{equation}

\begin{equation}
\varepsilon_{xz}=\frac{1}{2}\left(\frac{\partial u_x}{\partial z}+\frac{\partial u_z}{\partial x}\right).
\end{equation}

В термоупругой постановке полная деформация состоит из механической и тепловой частей:

\begin{equation}
\boldsymbol{\varepsilon}=\boldsymbol{\varepsilon}_{mech}+\boldsymbol{\varepsilon}_{th}.
\end{equation}

Тепловая деформация определяется выражением:

\begin{equation}
\boldsymbol{\varepsilon}_{th}=\alpha(T-T_{ref})\mathbf{I},
\end{equation}

где $\alpha$ --- коэффициент линейного теплового расширения, $T$ --- текущая температура, $T_{ref}$ --- опорная температура, $\mathbf{I}$ --- единичный тензор.

Механическая деформация:

\begin{equation}
\boldsymbol{\varepsilon}_{mech}=\boldsymbol{\varepsilon}-\boldsymbol{\varepsilon}_{th}.
\end{equation}

Напряжения определяются через закон Гука:

\begin{equation}
\boldsymbol{\sigma}=\mathbf{C}:\left(\boldsymbol{\varepsilon}-\boldsymbol{\varepsilon}_{th}\right),
\end{equation}

где $\mathbf{C}$ --- тензор упругости.

Для изотропного материала напряжения можно записать через параметры Ламе:

\begin{equation}
\boldsymbol{\sigma}=2\mu(\boldsymbol{\varepsilon}-\boldsymbol{\varepsilon}_{th})+\lambda\operatorname{tr}(\boldsymbol{\varepsilon}-\boldsymbol{\varepsilon}_{th})\mathbf{I}.
\end{equation}

Параметры Ламе выражаются через модуль Юнга $E$ и коэффициент Пуассона $\nu$:

\begin{equation}
\mu=\frac{E}{2(1+\nu)},
\end{equation}

\begin{equation}
\lambda=\frac{E\nu}{(1+\nu)(1-2\nu)}.
\end{equation}

Если расчёт выполняется в динамической постановке, уравнение движения имеет вид:

\begin{equation}
\rho\frac{\partial^2\mathbf{u}}{\partial t^2}=\nabla\cdot\boldsymbol{\sigma}+\mathbf{f}.
\end{equation}

Здесь $\rho$ --- плотность, а $\mathbf{f}$ --- объёмные силы. Динамическая постановка позволяет учитывать распространение механического возмущения во времени.

\section{Связанная термоупругая постановка в COMSOL}

В COMSOL задача реализована как связанная термоупругая модель. Это означает, что расчёт температуры и расчёт механического состояния не являются полностью независимыми. Температурное поле, полученное в тепловой части, используется в механической части для вычисления теплового расширения.

Связанная система может быть представлена следующим образом:

\begin{equation}
\rho C_p\frac{\partial T}{\partial t}=\nabla\cdot(k\nabla T)+Q,
\end{equation}

\begin{equation}
\rho\frac{\partial^2\mathbf{u}}{\partial t^2}=\nabla\cdot\boldsymbol{\sigma}+\mathbf{f},
\end{equation}

\begin{equation}
\boldsymbol{\sigma}=\mathbf{C}:\left[\boldsymbol{\varepsilon}(\mathbf{u})-\alpha(T-T_{ref})\mathbf{I}\right].
\end{equation}

В COMSOL данная постановка реализуется через следующие физические интерфейсы:

\begin{itemize}
    \item \textit{Heat Transfer in Solids} --- расчёт температурного поля;
    \item \textit{Solid Mechanics} --- расчёт перемещений, деформаций и напряжений;
    \item \textit{Thermal Expansion} --- связь между температурой и механической деформацией.
\end{itemize}

Интерфейс \textit{Heat Transfer in Solids} рассчитывает температуру в стержне и породе. Интерфейс \textit{Solid Mechanics} рассчитывает механический отклик. Связь \textit{Thermal Expansion} передаёт температурное поле в механическую часть и создаёт тепловую деформацию.

\section{Геометрическая модель}

Геометрия модели состоит из двух основных объектов:

\begin{enumerate}
    \item геологическая среда;
    \item металлический стержень.
\end{enumerate}

Геологическая среда задаётся как трёхмерный объём. В базовой постановке она может быть представлена кубом или прямоугольным параллелепипедом. Такая форма удобна для численного расчёта и анализа распространения возмущений в трёхмерном пространстве.

Металлический стержень задаётся как цилиндрический объект. Он имеет радиус $r_s$, длину $l_s$ и определённое положение относительно геологической среды. Стержень используется как нагретое тело, от которого начинается распространение тепла.

Объём стержня:

\begin{equation}
V_s=\pi r_s^2l_s.
\end{equation}

Площадь торца стержня:

\begin{equation}
S_{face}=\pi r_s^2.
\end{equation}

Площадь боковой поверхности:

\begin{equation}
S_{side}=2\pi r_sl_s.
\end{equation}

Если стержень контактирует с породой торцом, контактная площадь равна:

\begin{equation}
S_{contact}=\pi r_s^2.
\end{equation}

Если стержень частично внедрён в геологическую среду, контактная площадь может быть записана как:

\begin{equation}
S_{contact}=\pi r_s^2+2\pi r_sl_{in},
\end{equation}

где $l_{in}$ --- длина части стержня, находящейся внутри породы.

При построении геометрии важно обеспечить корректное взаимодействие стержня и породы. Для этого в COMSOL может использоваться режим \textit{Form Union}, при котором объекты объединяются в единую геометрию с общей внутренней границей. Такая постановка удобна для идеального теплового и механического контакта между стержнем и породой.

\section{Материалы и параметры модели}

После создания геометрии каждому домену назначаются физические свойства материала. В модели используются как минимум два типа материалов:

\begin{itemize}
    \item материал металлического стержня;
    \item материал геологической среды.
\end{itemize}

Для стержня задаются параметры, представленные в таблице~\ref{tab:rod_material_properties}.

\begin{table}[h!]
\centering
\caption{Основные свойства материала металлического стержня}
\label{tab:rod_material_properties}
\begin{tabular}{|l|c|c|p{6cm}|}
\hline
\textbf{Параметр} & \textbf{Обозначение} & \textbf{Единица} & \textbf{Назначение} \\
\hline
Плотность & $\rho_s$ & кг/м$^3$ & Используется в тепловой и механической частях \\
\hline
Теплопроводность & $k_s$ & Вт/(м$\cdot$К) & Определяет скорость переноса тепла \\
\hline
Удельная теплоёмкость & $C_{p,s}$ & Дж/(кг$\cdot$К) & Определяет тепловую инерционность \\
\hline
Модуль Юнга & $E_s$ & Па & Определяет жёсткость материала \\
\hline
Коэффициент Пуассона & $\nu_s$ & 1 & Описывает поперечную деформацию \\
\hline
Коэффициент теплового расширения & $\alpha_s$ & 1/К & Определяет расширение при нагреве \\
\hline
\end{tabular}
\end{table}

Для геологической среды задаются параметры, представленные в таблице~\ref{tab:rock_material_properties}.

\begin{table}[h!]
\centering
\caption{Основные свойства материала геологической среды}
\label{tab:rock_material_properties}
\begin{tabular}{|l|c|c|p{6cm}|}
\hline
\textbf{Параметр} & \textbf{Обозначение} & \textbf{Единица} & \textbf{Назначение} \\
\hline
Плотность & $\rho_r$ & кг/м$^3$ & Масса единицы объёма породы \\
\hline
Теплопроводность & $k_r$ & Вт/(м$\cdot$К) & Способность породы проводить тепло \\
\hline
Удельная теплоёмкость & $C_{p,r}$ & Дж/(кг$\cdot$К) & Энергия, необходимая для нагрева \\
\hline
Модуль Юнга & $E_r$ & Па & Упругая жёсткость породы \\
\hline
Коэффициент Пуассона & $\nu_r$ & 1 & Механическая характеристика деформации \\
\hline
Коэффициент теплового расширения & $\alpha_r$ & 1/К & Тепловое расширение породы \\
\hline
\end{tabular}
\end{table}

Эти параметры используются COMSOL при решении тепловой и механической частей задачи. Теплопроводность, плотность и теплоёмкость входят в уравнение теплопроводности. Модуль Юнга, коэффициент Пуассона и плотность используются в механике твёрдого тела. Коэффициент теплового расширения используется в связи \textit{Thermal Expansion}.

\section{Начальные условия}

Начальные условия задают состояние модели в момент времени $t=0$. В рассматриваемой симуляции основное начальное воздействие связано с разностью температур между металлическим стержнем и геологической средой.

Для геологической среды задаётся начальная температура:

\begin{equation}
T_{rock}(x,y,z,0)=T_0.
\end{equation}

Для металлического стержня задаётся повышенная начальная температура:

\begin{equation}
T_{rod}(x,y,z,0)=T_s.
\end{equation}

При этом выполняется условие:

\begin{equation}
T_s>T_0.
\end{equation}

Разность температур:

\begin{equation}
\Delta T=T_s-T_0
\end{equation}

создаёт начальный температурный градиент и запускает процесс теплопередачи.

Для механической части начальные перемещения принимаются равными нулю:

\begin{equation}
\mathbf{u}(x,y,z,0)=\mathbf{0}.
\end{equation}

Если в расчёте учитывается динамика, начальные скорости также задаются равными нулю:

\begin{equation}
\frac{\partial\mathbf{u}}{\partial t}(x,y,z,0)=\mathbf{0}.
\end{equation}

Это означает, что до начала теплового воздействия среда находится в исходном механическом состоянии.

\section{Граничные условия}

Граничные условия определяют поведение модели на внешних поверхностях и в области контакта между стержнем и породой.

Для тепловой части могут использоваться условия теплоизоляции, конвективного теплообмена или заданной температуры.

Условие теплоизоляции:

\begin{equation}
-k\nabla T\cdot\mathbf{n}=0.
\end{equation}

Оно означает, что через границу не проходит тепловой поток.

Условие конвективного теплообмена:

\begin{equation}
-k\nabla T\cdot\mathbf{n}=h(T-T_{amb}),
\end{equation}

где $h$ --- коэффициент теплоотдачи, а $T_{amb}$ --- температура окружающей среды.

Условие заданной температуры:

\begin{equation}
T=T_b.
\end{equation}

Для механической части используются условия закрепления и свободных поверхностей.

Полное закрепление:

\begin{equation}
\mathbf{u}=\mathbf{0}.
\end{equation}

Частичное закрепление, например по оси $Z$:

\begin{equation}
u_z=0.
\end{equation}

Свободная поверхность:

\begin{equation}
\boldsymbol{\sigma}\mathbf{n}=\mathbf{0}.
\end{equation}

В области контакта стержня и породы задаётся передача тепла и механического воздействия. Для тепловой части это означает согласование температуры и теплового потока на общей границе.

\section{Конечно-элементная сетка}

COMSOL решает задачу методом конечных элементов. Для этого геометрия разбивается на конечное число малых элементов. Расчётная область $\Omega$ аппроксимируется как объединение элементов:

\begin{equation}
\Omega\approx\bigcup_{e=1}^{N_e}\Omega_e,
\end{equation}

где $N_e$ --- количество конечных элементов.

Узлы сетки обозначаются как:

\begin{equation}
\mathcal{V}=\{v_1,v_2,...,v_N\},
\end{equation}

где $N$ --- количество узлов.

Каждый узел имеет координаты:

\begin{equation}
\mathbf{x}_i=[x_i,y_i,z_i]^T.
\end{equation}

Температура аппроксимируется через функции формы:

\begin{equation}
T(\mathbf{x},t)\approx\sum_{i=1}^{N}N_i(\mathbf{x})T_i(t).
\end{equation}

Поле перемещений аппроксимируется аналогично:

\begin{equation}
\mathbf{u}(\mathbf{x},t)\approx\sum_{i=1}^{N}N_i(\mathbf{x})\mathbf{u}_i(t).
\end{equation}

Для данной задачи особое значение имеет качество сетки в зоне контакта стержня и породы. Именно там возникают наибольшие температурные градиенты, деформации и напряжения. Поэтому в контактной области желательно использовать более мелкую сетку, а в удалённых областях можно использовать более крупные элементы.

Сетка также важна для последующего экспорта данных. Координаты узлов, элементы и связность сетки должны быть сохранены, чтобы расчётные поля можно было корректно интерпретировать как пространственные данные.

\section{Временной расчёт}

Симуляция выполняется во времени, так как процесс теплопередачи и термоупругого отклика является нестационарным. Расчётный интервал задаётся как:

\begin{equation}
t\in[0,T_{end}].
\end{equation}

Он состоит из временных точек:

\begin{equation}
t_0,t_1,t_2,...,t_M.
\end{equation}

Если шаг сохранения постоянный:

\begin{equation}
t_{m+1}=t_m+\Delta t.
\end{equation}

В COMSOL необходимо различать внутренний шаг решателя и шаг экспорта данных. Решатель может использовать адаптивный шаг для устойчивого решения уравнений, а экспорт результатов может выполняться через фиксированный временной интервал. Для дальнейшей обработки данных фиксированный шаг экспорта является удобным, так как все физические поля сохраняются в виде последовательности состояний.

На каждом временном шаге COMSOL рассчитывает температуру, перемещения, скорости, деформации и напряжения в узлах расчётной области. Таким образом, результатом расчёта является пространственно-временной набор данных.

\section{Расчётные физические поля}

В результате COMSOL-симуляции формируются несколько групп физических величин.

\subsection{Температурное поле}

Температурное поле:

\begin{equation}
T=T(x,y,z,t)
\end{equation}

показывает распределение температуры в стержне и геологической среде. Для каждого узла и временного шага сохраняется значение температуры:

\begin{equation}
T_i^t.
\end{equation}

Температура является основной величиной тепловой части расчёта.

\subsection{Тепловой поток}

Тепловой поток определяется выражением:

\begin{equation}
\mathbf{q}=-k\nabla T.
\end{equation}

Он показывает направление и интенсивность переноса тепла. Компоненты теплового потока:

\begin{equation}
\mathbf{q}=[q_x,q_y,q_z].
\end{equation}

Модуль теплового потока:

\begin{equation}
|\mathbf{q}|=\sqrt{q_x^2+q_y^2+q_z^2}.
\end{equation}

\subsection{Перемещения и скорости}

Поле перемещений:

\begin{equation}
\mathbf{u}=[u_x,u_y,u_z]
\end{equation}

показывает смещение материала из исходного положения.

Модуль перемещения:

\begin{equation}
|\mathbf{u}|=\sqrt{u_x^2+u_y^2+u_z^2}.
\end{equation}

Скорости перемещений:

\begin{equation}
\dot{\mathbf{u}}=[\dot{u}_x,\dot{u}_y,\dot{u}_z]
\end{equation}

показывают изменение перемещений во времени.

\subsection{Деформации}

Деформации показывают относительное изменение формы и размеров материала. Они являются безразмерными величинами:

\begin{equation}
\varepsilon=\frac{\Delta L}{L}.
\end{equation}

В трёхмерной задаче деформационное состояние описывается тензором:

\begin{equation}
\boldsymbol{\varepsilon}=\begin{bmatrix}
\varepsilon_{xx} & \varepsilon_{xy} & \varepsilon_{xz} \\
\varepsilon_{yx} & \varepsilon_{yy} & \varepsilon_{yz} \\
\varepsilon_{zx} & \varepsilon_{zy} & \varepsilon_{zz}
\end{bmatrix}.
\end{equation}

\subsection{Напряжения}

Напряжения показывают внутреннюю механическую нагрузку материала. Тензор напряжений:

\begin{equation}
\boldsymbol{\sigma}=\begin{bmatrix}
\sigma_{xx} & \sigma_{xy} & \sigma_{xz} \\
\sigma_{yx} & \sigma_{yy} & \sigma_{yz} \\
\sigma_{zx} & \sigma_{zy} & \sigma_{zz}
\end{bmatrix}.
\end{equation}

Для симметричного тензора:

\begin{equation}
\sigma_{xy}=\sigma_{yx}, \quad
\sigma_{yz}=\sigma_{zy}, \quad
\sigma_{xz}=\sigma_{zx}.
\end{equation}

Также может использоваться эквивалентное напряжение по Мизесу:

\begin{equation}
\sigma_{vm}=\sqrt{\frac{1}{2}\left[(\sigma_{xx}-\sigma_{yy})^2+(\sigma_{yy}-\sigma_{zz})^2+(\sigma_{zz}-\sigma_{xx})^2+6(\sigma_{xy}^2+\sigma_{yz}^2+\sigma_{xz}^2)\right]}.
\end{equation}

\section{Экспорт данных из COMSOL}

После выполнения расчёта данные экспортируются из COMSOL в набор файлов. Экспорт выполняется из набора данных \texttt{Study 1/Solution 1 (sol1)} с параметром \texttt{Time selection: All}. Это означает, что в файлы сохраняются значения по всем рассчитанным временным шагам.

В текущей структуре одна симуляция содержит следующие файлы:

\begin{table}[h!]
\centering
\caption{Файлы, полученные при экспорте одной COMSOL-симуляции}
\label{tab:exported_files}
\begin{tabular}{|l|p{9cm}|}
\hline
\textbf{Файл} & \textbf{Содержание} \\
\hline
\texttt{data\_temperature.csv} & Температура, координаты и тепловые свойства \\
\hline
\texttt{data\_displacement.csv} & Перемещения и скорости перемещений \\
\hline
\texttt{data\_materials.csv} & Механические и термоупругие свойства материалов \\
\hline
\texttt{data\_strain.csv} & Компоненты деформаций \\
\hline
\texttt{data\_stress\_1.csv} & Основные компоненты напряжений \\
\hline
\texttt{data\_stress\_2.csv} & Сдвиговые компоненты напряжений \\
\hline
\texttt{data\_stress\_3.csv} & Дополнительные компоненты деформации \\
\hline
\texttt{sandstone.mphtxt} & Конечно-элементная сетка симуляции \\
\hline
\end{tabular}
\end{table}

Такое разделение позволяет отдельно хранить разные группы физических величин и затем объединять их на этапе обработки данных.

\section{Описание файла \texttt{data\_temperature.csv}}

Файл \texttt{data\_temperature.csv} содержит данные тепловой части расчёта.

\begin{table}[h!]
\centering
\caption{Поля файла \texttt{data\_temperature.csv}}
\label{tab:data_temperature}
\begin{tabular}{|l|c|p{8cm}|}
\hline
\textbf{Поле} & \textbf{Единица} & \textbf{Описание} \\
\hline
\texttt{T} & K & Температура в точке расчётной области \\
\hline
\texttt{x} & m & Координата по оси X \\
\hline
\texttt{y} & m & Координата по оси Y \\
\hline
\texttt{z} & m & Координата по оси Z \\
\hline
\texttt{ht.k\_iso} & W/(m$\cdot$K) & Изотропная теплопроводность \\
\hline
\texttt{ht.rho} & kg/m$^3$ & Плотность в тепловом интерфейсе \\
\hline
\texttt{ht.Cp} & J/(kg$\cdot$K) & Удельная теплоёмкость \\
\hline
\end{tabular}
\end{table}

Температура \texttt{T} показывает тепловое состояние среды в заданной точке и момент времени. Координаты \texttt{x}, \texttt{y}, \texttt{z} позволяют определить пространственное положение этой точки.

Параметр \texttt{ht.k\_iso} отвечает за теплопроводность материала. Параметр \texttt{ht.rho} задаёт плотность, а \texttt{ht.Cp} --- теплоёмкость. Вместе они определяют тепловую инерционность материала:

\begin{equation}
\rho C_p.
\end{equation}

Этот файл является основным источником данных о температурном поле модели.

\section{Описание файла \texttt{data\_displacement.csv}}

Файл \texttt{data\_displacement.csv} содержит данные о перемещениях и скоростях.

\begin{table}[h!]
\centering
\caption{Поля файла \texttt{data\_displacement.csv}}
\label{tab:data_displacement}
\begin{tabular}{|l|c|p{8cm}|}
\hline
\textbf{Поле} & \textbf{Единица} & \textbf{Описание} \\
\hline
\texttt{u} & m & Перемещение по оси X \\
\hline
\texttt{v} & m & Перемещение по оси Y \\
\hline
\texttt{w} & m & Перемещение по оси Z \\
\hline
\texttt{ut} & m/s & Скорость перемещения по оси X \\
\hline
\texttt{vt} & m/s & Скорость перемещения по оси Y \\
\hline
\texttt{wt} & m/s & Скорость перемещения по оси Z \\
\hline
\end{tabular}
\end{table}

В COMSOL компоненты \texttt{u}, \texttt{v}, \texttt{w} соответствуют перемещениям по координатным направлениям:

\begin{equation}
u=u_x, \quad v=u_y, \quad w=u_z.
\end{equation}

Вектор перемещения:

\begin{equation}
\mathbf{u}_i^t=[u_i^t,v_i^t,w_i^t]^T.
\end{equation}

Модуль перемещения:

\begin{equation}
|\mathbf{u}_i^t|=\sqrt{(u_i^t)^2+(v_i^t)^2+(w_i^t)^2}.
\end{equation}

Компоненты \texttt{ut}, \texttt{vt}, \texttt{wt} показывают скорость изменения перемещений во времени:

\begin{equation}
\dot{\mathbf{u}}_i^t=[ut_i^t,vt_i^t,wt_i^t]^T.
\end{equation}

Данный файл описывает механическую реакцию среды на тепловое воздействие.

\section{Описание файла \texttt{data\_materials.csv}}

Файл \texttt{data\_materials.csv} содержит свойства материалов, использованные в механической и термоупругой частях модели.

\begin{table}[h!]
\centering
\caption{Поля файла \texttt{data\_materials.csv}}
\label{tab:data_materials}
\begin{tabular}{|l|c|p{8cm}|}
\hline
\textbf{Поле} & \textbf{Единица} & \textbf{Описание} \\
\hline
\texttt{solid.E} & Pa & Модуль Юнга \\
\hline
\texttt{solid.nu} & 1 & Коэффициент Пуассона \\
\hline
\texttt{solid.rho} & kg/m$^3$ & Плотность в механическом интерфейсе \\
\hline
\texttt{te1.alpha\_iso} & 1/K & Изотропный коэффициент теплового расширения \\
\hline
\end{tabular}
\end{table}

Модуль Юнга \texttt{solid.E} определяет жёсткость материала. Коэффициент Пуассона \texttt{solid.nu} описывает поперечную деформацию при продольном воздействии. Плотность \texttt{solid.rho} используется в механической части расчёта. Коэффициент \texttt{te1.alpha\_iso} определяет тепловое расширение материала.

Тепловая деформация рассчитывается как:

\begin{equation}
\boldsymbol{\varepsilon}_{th}=\alpha(T-T_{ref})\mathbf{I}.
\end{equation}

Этот файл фиксирует физические свойства материалов, необходимые для интерпретации результатов симуляции.

\section{Описание файла \texttt{data\_strain.csv}}

Файл \texttt{data\_strain.csv} содержит компоненты деформационного состояния материала. Деформация показывает относительное изменение формы или размеров тела.

Общий вид тензора деформаций:

\begin{equation}
\boldsymbol{\varepsilon}=\begin{bmatrix}
\varepsilon_{xx} & \varepsilon_{xy} & \varepsilon_{xz} \\
\varepsilon_{yx} & \varepsilon_{yy} & \varepsilon_{yz} \\
\varepsilon_{zx} & \varepsilon_{zy} & \varepsilon_{zz}
\end{bmatrix}.
\end{equation}

Деформации являются безразмерными величинами. Они могут быть использованы для анализа областей растяжения, сжатия и сдвига в геологической среде.

Если в файле экспортированы отдельные компоненты деформации, каждая из них описывает изменение материала в соответствующем направлении или плоскости.

\section{Описание файла \texttt{data\_stress\_1.csv}}

Файл \texttt{data\_stress\_1.csv} используется для хранения основных компонент напряжений. В типовой структуре сюда относятся нормальные компоненты тензора напряжений.

\begin{table}[h!]
\centering
\caption{Поля файла \texttt{data\_stress\_1.csv}}
\label{tab:data_stress_1}
\begin{tabular}{|l|c|p{8cm}|}
\hline
\textbf{Поле} & \textbf{Единица} & \textbf{Описание} \\
\hline
\texttt{solid.sx} & N/m$^2$ & Нормальное напряжение по оси X \\
\hline
\texttt{solid.sy} & N/m$^2$ & Нормальное напряжение по оси Y \\
\hline
\texttt{solid.sz} & N/m$^2$ & Нормальное напряжение по оси Z \\
\hline
\end{tabular}
\end{table}

Нормальные напряжения показывают растяжение или сжатие материала вдоль соответствующих направлений. Они являются диагональными компонентами тензора напряжений:

\begin{equation}
\sigma_{xx}, \quad \sigma_{yy}, \quad \sigma_{zz}.
\end{equation}

Этот файл описывает основную часть напряжённого состояния среды.

\section{Описание файла \texttt{data\_stress\_2.csv}}

Файл \texttt{data\_stress\_2.csv} содержит сдвиговые компоненты напряжений.

\begin{table}[h!]
\centering
\caption{Поля файла \texttt{data\_stress\_2.csv}}
\label{tab:data_stress_2}
\begin{tabular}{|l|c|p{8cm}|}
\hline
\textbf{Поле} & \textbf{Единица} & \textbf{Описание} \\
\hline
\texttt{solid.sxy} & N/m$^2$ & Сдвиговое напряжение в плоскости XY \\
\hline
\texttt{solid.syz} & N/m$^2$ & Сдвиговое напряжение в плоскости YZ \\
\hline
\texttt{solid.sxz} & N/m$^2$ & Сдвиговое напряжение в плоскости XZ \\
\hline
\end{tabular}
\end{table}

Сдвиговые напряжения описывают касательные нагрузки в материале. Совместно с нормальными компонентами они позволяют восстановить тензор напряжений:

\begin{equation}
\boldsymbol{\sigma}=\begin{bmatrix}
\texttt{solid.sx} & \texttt{solid.sxy} & \texttt{solid.sxz} \\
\texttt{solid.sxy} & \texttt{solid.sy} & \texttt{solid.syz} \\
\texttt{solid.sxz} & \texttt{solid.syz} & \texttt{solid.sz}
\end{bmatrix}.
\end{equation}

\section{Описание файла \texttt{data\_stress\_3.csv}}

Файл \texttt{data\_stress\_3.csv} в текущей структуре содержит поля, представленные в таблице~\ref{tab:data_stress_3}.

\begin{table}[h!]
\centering
\caption{Поля файла \texttt{data\_stress\_3.csv}}
\label{tab:data_stress_3}
\begin{tabular}{|l|c|p{8cm}|}
\hline
\textbf{Поле} & \textbf{Единица} & \textbf{Описание} \\
\hline
\texttt{solid.eX} & 1 & Компонента деформации по направлению X \\
\hline
\texttt{solid.eY} & 1 & Компонента деформации по направлению Y \\
\hline
\texttt{solid.eZ} & 1 & Компонента деформации по направлению Z \\
\hline
\end{tabular}
\end{table}

Несмотря на название файла, данные поля относятся к деформациям, а не к напряжениям. Это видно по обозначению \texttt{solid.eX}, \texttt{solid.eY}, \texttt{solid.eZ} и по единице измерения $1$, так как деформация является безразмерной величиной.

Данный файл сохраняется в общей структуре датасета как дополнительный файл механических характеристик, содержащий деформационные компоненты.

\section{Описание файла \texttt{sandstone.mphtxt}}

Файл \texttt{sandstone.mphtxt} является текстовым файлом COMSOL, содержащим конечно-элементную сетку симуляции. Он отличается от CSV-файлов тем, что хранит не физические поля, а структуру mesh-модели.

Данный файл содержит информацию, необходимую для восстановления расчётной сетки:

\begin{itemize}
    \item координаты узлов;
    \item состав конечных элементов;
    \item связность между узлами;
    \item принадлежность элементов к доменам;
    \item топологию конечно-элементной модели.
\end{itemize}

Файл \texttt{sandstone.mphtxt} является важной частью набора данных, потому что физические значения из CSV-файлов имеют смысл только вместе с пространственной структурой, на которой они были рассчитаны.

Если данные используются в графовом представлении, узлы конечно-элементной сетки становятся узлами графа, а связи внутри элементов становятся рёбрами. В таком случае \texttt{sandstone.mphtxt} используется для построения графовой структуры:

\begin{equation}
\mathcal{G}=(\mathcal{V},\mathcal{E}),
\end{equation}

где $\mathcal{V}$ --- множество узлов сетки, а $\mathcal{E}$ --- множество связей между ними.

\section{Объединение экспортированных данных}

Все экспортированные файлы относятся к одной симуляции и описывают разные физические аспекты одного процесса. Для дальнейшей обработки данные могут быть объединены по временным шагам и пространственным координатам.

Основными идентификаторами являются:

\begin{equation}
time,
\end{equation}

\begin{equation}
x,y,z.
\end{equation}

Если на этапе предобработки формируется явный идентификатор узла, можно использовать:

\begin{equation}
node\_id.
\end{equation}

Итоговая строка объединённого датасета может соответствовать одной точке расчётной области в один момент времени и включать поля, представленные в таблице~\ref{tab:combined_dataset}.

\begin{table}[h!]
\centering
\caption{Пример структуры объединённой строки датасета}
\label{tab:combined_dataset}
\begin{tabular}{|p{4cm}|p{9cm}|}
\hline
\textbf{Группа полей} & \textbf{Поля} \\
\hline
Идентификаторы & \texttt{time}, \texttt{node\_id} \\
\hline
Координаты & \texttt{x}, \texttt{y}, \texttt{z} \\
\hline
Температура & \texttt{T} \\
\hline
Перемещения & \texttt{u}, \texttt{v}, \texttt{w} \\
\hline
Скорости & \texttt{ut}, \texttt{vt}, \texttt{wt} \\
\hline
Напряжения & \texttt{sx}, \texttt{sy}, \texttt{sz}, \texttt{sxy}, \texttt{syz}, \texttt{sxz} \\
\hline
Деформации & \texttt{eX}, \texttt{eY}, \texttt{eZ} \\
\hline
Тепловые параметры & \texttt{k\_iso}, \texttt{rho}, \texttt{Cp} \\
\hline
Механические параметры & \texttt{E}, \texttt{nu}, \texttt{alpha} \\
\hline
\end{tabular}
\end{table}

Такой формат позволяет хранить полное физическое состояние симуляции в табличном виде. При этом файл \texttt{sandstone.mphtxt} хранится отдельно, так как он описывает сеточную структуру, а не значения физических полей.

\section{Структура папки одной симуляции}

Для хранения данных одной симуляции используется отдельная папка. Пример структуры:

\begin{verbatim}
simulation_001/
  data_temperature.csv
  data_displacement.csv
  data_materials.csv
  data_strain.csv
  data_stress_1.csv
  data_stress_2.csv
  data_stress_3.csv
  sandstone.mphtxt
\end{verbatim}

Такая структура позволяет хранить все данные одной симуляции вместе. При увеличении количества симуляций можно использовать нумерацию:

\begin{verbatim}
dataset/
  sim_0001/
    data_temperature.csv
    data_displacement.csv
    data_materials.csv
    data_strain.csv
    data_stress_1.csv
    data_stress_2.csv
    data_stress_3.csv
    stone001.mphtxt
  sim_0002/
    data_temperature.csv
    data_displacement.csv
    data_materials.csv
    data_strain.csv
    data_stress_1.csv
    data_stress_2.csv
    data_stress_3.csv
    stone002.mphtxt
\end{verbatim}

Каждая папка соответствует одному расчётному сценарию COMSOL.

\section{Практическая организация экспорта в COMSOL}

Экспорт расчётных данных в COMSOL выполняется после завершения временного расчёта. В данной работе для каждого типа физических величин создаётся отдельный узел экспорта. Такой подход позволяет не смешивать разные группы данных в одном файле и делает структуру набора данных более понятной.

Каждый экспортируемый файл создаётся из решения:

\begin{equation}
\texttt{Study 1/Solution 1 (sol1)}.
\end{equation}

Выбор данного набора данных означает, что экспорт выполняется по результатам первой расчётной задачи и первого решения модели. Для всех файлов используется параметр:

\begin{equation}
\texttt{Time selection: All}.
\end{equation}

Это означает, что COMSOL сохраняет значения выбранных выражений для всех временных шагов, доступных в решении. Таким образом, один CSV-файл содержит не статическое распределение величины, а полную временную историю соответствующего поля.

Например, если симуляция рассчитана в моменты времени:

\begin{equation}
t_0,t_1,t_2,\ldots,t_M,
\end{equation}

то файл \texttt{data\_temperature.csv} содержит температуру не только при $t_0$, но и для всех последующих моментов времени:

\begin{equation}
T^0,T^1,T^2,\ldots,T^M.
\end{equation}

Аналогично файл \texttt{data\_displacement.csv} содержит перемещения и скорости на всех временных шагах, а файлы напряжений и деформаций содержат соответствующие механические поля во времени.

Практически это означает, что один файл описывает поле в виде пространственно-временной таблицы. Каждая строка относится к некоторой точке расчётной области и некоторому моменту времени. Поэтому при дальнейшей обработке необходимо учитывать не только координаты точки, но и временной слой, к которому относится значение.

\section{Логика разделения экспортируемых данных по файлам}

Разделение экспорта на несколько файлов связано с тем, что физические величины имеют разную природу и разные единицы измерения. Температура, перемещения, напряжения, деформации и свойства материалов относятся к разным группам данных.

Температурный файл содержит величины тепловой задачи. Он описывает распределение температуры и параметры, участвующие в уравнении теплопроводности. Файл перемещений содержит механическую реакцию среды. Файлы напряжений и деформаций описывают внутреннее механическое состояние материала. Файл материалов содержит параметры, которые определяют поведение доменов в тепловой и механической частях.

Такое разделение удобно по нескольким причинам:

\begin{enumerate}
    \item упрощается контроль содержимого каждого файла, так как каждая таблица отвечает за конкретную физическую группу;
    \item снижается риск перепутать величины с разными единицами измерения;
    \item можно отдельно проверять корректность температуры, перемещений, напряжений и деформаций;
    \item при последующей обработке можно выбирать только те файлы, которые нужны для конкретной задачи;
    \item структура становится расширяемой: при необходимости можно добавить новый файл с дополнительными физическими полями.
\end{enumerate}

В результате каждый файл является отдельным логическим блоком общей симуляции, а полное состояние расчёта восстанавливается через совместное использование всех экспортированных данных.

\section{Связь между файлами одной симуляции}

Все CSV-файлы, экспортированные из одной симуляции, относятся к одной и той же геометрии, одной и той же сетке и одному и тому же временному расчёту. Поэтому они должны рассматриваться как части единого набора данных.

Связь между файлами осуществляется через пространственное положение и время. В идеальном случае каждая строка может быть идентифицирована парой:

\begin{equation}
(time,node\_id),
\end{equation}

где \texttt{time} --- временной шаг, а \texttt{node\_id} --- номер узла конечно-элементной сетки.

Если явный номер узла не выгружается напрямую, соответствие между строками может быть восстановлено через координаты:

\begin{equation}
(time,x,y,z).
\end{equation}

При этом важно, чтобы координаты во всех файлах соответствовали одной и той же сетке. Если разные физические поля экспортируются на разных наборах точек, объединение данных усложняется. Поэтому в рамках одной симуляции желательно использовать одинаковый набор точек экспорта для всех физических величин.

Полное состояние одного узла в момент времени $t$ может быть представлено как вектор:

\begin{equation}
\mathbf{s}_i^t =
\left[
x_i,
y_i,
z_i,
T_i^t,
u_i^t,
v_i^t,
w_i^t,
ut_i^t,
vt_i^t,
wt_i^t,
\sigma_i^t,
\varepsilon_i^t,
\mathbf{m}_i
\right],
\end{equation}

где $\sigma_i^t$ обозначает набор компонент напряжений, $\varepsilon_i^t$ --- набор компонент деформаций, а $\mathbf{m}_i$ --- свойства материала в данной точке.

Такое представление удобно, потому что оно объединяет тепловые, механические и материальные характеристики в одну структуру данных.

\section{Описание временной структуры экспортированных данных}

Так как в настройках экспорта выбран параметр \texttt{Time selection: All}, данные имеют временную структуру. Пусть расчёт содержит $M+1$ временных состояний:

\begin{equation}
t_0,t_1,\ldots,t_M.
\end{equation}

Для каждого временного состояния сохраняется распределение физических величин по расчётной области. Если сетка содержит $N$ узлов, то для одного поля формируется набор значений:

\begin{equation}
\left\{y_i^t\right\}_{i=1}^{N},
\end{equation}

где $y_i^t$ --- значение физической величины в узле $i$ в момент времени $t$.

Для всех временных шагов это можно представить как матрицу:

\begin{equation}
\mathbf{Y}=
\begin{bmatrix}
y_1^{t_0} & y_2^{t_0} & \ldots & y_N^{t_0} \\
y_1^{t_1} & y_2^{t_1} & \ldots & y_N^{t_1} \\
\vdots & \vdots & \ddots & \vdots \\
y_1^{t_M} & y_2^{t_M} & \ldots & y_N^{t_M}
\end{bmatrix}.
\end{equation}

Для температуры такая матрица описывает эволюцию температурного поля. Для перемещений аналогичная структура строится отдельно для компонент $u$, $v$ и $w$. Для напряжений и деформаций структура строится для каждой компоненты соответствующего тензора.

Временная структура особенно важна, потому что рассматриваемый процесс является динамическим. Значения физических полей на поздних шагах зависят от предыдущих состояний системы. Поэтому при обработке данных нельзя рассматривать временные строки как полностью независимые наблюдения.

\section{Описание пространственной структуры экспортированных данных}

Пространственная структура данных задаётся координатами точек и конечно-элементной сеткой. Координаты $x$, $y$, $z$, экспортируемые вместе с температурным полем, определяют положение точек в трёхмерном пространстве:

\begin{equation}
\mathbf{x}_i=[x_i,y_i,z_i]^T.
\end{equation}

Однако сами координаты описывают только положение точек. Для восстановления полной структуры расчётной области необходима информация о том, какие точки соединены между собой элементами сетки. Эту информацию содержит файл \texttt{sandstone.mphtxt}.

Конечно-элементная сетка может быть представлена как множество узлов и элементов:

\begin{equation}
\mathcal{M}=(\mathcal{V},\mathcal{K}),
\end{equation}

где $\mathcal{V}$ --- множество узлов, а $\mathcal{K}$ --- множество конечных элементов.

Если элемент является тетраэдром, он может быть задан четырьмя узлами:

\begin{equation}
K_e=\{v_a,v_b,v_c,v_d\}.
\end{equation}

Наличие такой информации позволяет восстановить локальное соседство точек. Это важно для интерпретации распространения тепла и механических возмущений, поскольку физические процессы в сплошной среде зависят от локальных взаимодействий между соседними областями.

\section{Контроль корректности COMSOL-симуляции}

Для того чтобы результаты симуляции были пригодны для дальнейшей работы, необходимо контролировать корректность модели. Проверка включает несколько аспектов.

Во-первых, проверяется геометрия: стержень должен быть расположен корректно относительно геологической среды, а область контакта должна быть сформирована правильно.

Во-вторых, проверяются материалы: каждому домену должны быть назначены соответствующие физические свойства. Ошибка в назначении материала может привести к полностью некорректному тепловому или механическому отклику.

В-третьих, проверяются начальные условия: стержень должен иметь повышенную начальную температуру, а порода --- исходную температуру.

В-четвёртых, проверяются граничные условия: они должны соответствовать физической постановке задачи.

В-пятых, проверяется сетка: в области контакта она должна быть достаточно детальной, а элементы не должны иметь плохое качество.

В-шестых, проверяются расчётные поля. Температура, перемещения, напряжения и деформации должны изменяться физически правдоподобно. Не должно быть неестественных скачков, бесконечных значений или явно ошибочных результатов.

\section{Итог по COMSOL-части}

В данной главе была описана COMSOL-часть общей дипломной работы. В ней была создана численная модель термоупругого процесса в геологической среде под воздействием нагретого металлического стержня.

В модели были использованы теплопередача в твёрдых телах, механика твёрдого тела и связь через тепловое расширение. Геологическая среда и стержень были представлены как отдельные домены с собственными физическими свойствами. На основе заданных начальных и граничных условий был выполнен временной расчёт, в результате которого получены температурные, механические и термоупругие поля.

Результаты симуляции были экспортированы в набор файлов: \texttt{data\_temperature.csv}, \texttt{data\_displacement.csv}, \texttt{data\_materials.csv}, \texttt{data\_strain.csv}, \texttt{data\_stress\_1.csv}, \texttt{data\_stress\_2.csv}, \texttt{data\_stress\_3.csv} и \texttt{sandstone.mphtxt}. CSV-файлы содержат численные значения физических величин, а файл \texttt{sandstone.mphtxt} хранит конечно-элементную сетку симуляции.

Таким образом, COMSOL-часть представляет собой документированную численную модель и набор расчётных данных, которые описывают развитие термоупругого процесса во времени и пространстве.
